\documentclass[10pt,a4paper]{amsart}
\usepackage[margin=19mm]{geometry}
\usepackage{amsmath,amssymb,mathtools}
\usepackage[scr=rsfs,cal=euler]{mathalfa}
\usepackage{xfrac}
\usepackage[unicode,hidelinks,bookmarks=false]{hyperref}
\newcommand{\ie}{i.e.\ }
\newcommand{\cf}{cf.\ }
\newcommand{\resp}{respectively\ }
\newcommand{\Subsection}{\subsection}
\providecommand{\nobreakdash}{\nobreak\hbox{-}}
\setlength{\emergencystretch}{2em}
\multlinegap0pt
\usepackage{bm}
\usepackage{bbm}
\usepackage{dsfont}
\usepackage{xspace}
\usepackage{cancel}
\usepackage{mathtools}
\usepackage{amsthm}
\makeatletter
\@ifundefined{Theorem}{\newtheorem{Theorem}{Theorem}}{}
\@ifundefined{Definition}{\newtheorem{Definition}[Theorem]{Definition}}{}
\@ifundefined{Notation}{\newtheorem{Notation}[Theorem]{Notation}}{}
\@ifundefined{Proposition}{\newtheorem{Proposition}[Theorem]{Proposition}}{}
\makeatother
\newcommand{\R}{\mathbb{R}}
\newcommand{\diag}{\operatorname{diag}}
\newcommand{\tr}{\operatorname{tr}}
\title[Density-Matrix Spectral Embeddings for Categorical Data: Ope]{Density-Matrix Spectral Embeddings for Categorical Data: Operator Structure and Stability}
\author{Raquel Bosch-Romeu, Antonio Falc\'o, os\'e-Antonio Rodr\'iguez-Gallego}
\date{Source: arXiv:2603.01975v1 (CC BY 4.0)}
\begin{document}
\maketitle

\noindent\emph{Source and licence.} Density-Matrix Spectral Embeddings for Categorical Data: Operator Structure and Stability. Version arXiv:2603.01975v1 (CC BY 4.0), distributed under the Creative Commons Attribution 4.0 International licence, as recorded in the arXiv licence metadata for this exact version and shown on the abstract page https://arxiv.org/abs/2603.01975v1. Full rights notice: licenses/arxiv-2603.01975-CC-BY-4.0.txt.\par\smallskip

\noindent\textit{Editorial scope.} Counted unit: the Proposition whose source label is \texttt{prop:exact\_relation\_rho\_hat\_rho} in the article (source0.tex lines 1064--1082, source label \texttt{prop:exact\_relation\_rho\_hat\_rho}), delivered together with the article's own complete original proof of it (source0.tex lines 1084--1109, one \texttt{proof} environment of 26 lines). No mathematics is written, repaired or re-derived: statement and proof are the article's own text line for line. Required context retained uncounted: def:class\_normalized\_profiles (lines 522--551); not:class\_masses\_profiles (lines 1043--1062). Every definition, display and label of those lines is retained unchanged. Omitted from this package and not redistributed: 1--521, 552--1042, 1063--1063, 1083--1083, 1110--1946. Nothing inside a retained excerpt is omitted or altered: every retained body is delivered exactly as the article's lines are written, so no omission receipt of any kind accompanies this package. Citations: the retained excerpts carry no citation command at all, so no bibliography entry of the article is transcribed and no reference is redistributed. Reference adaptation: none was needed and none was made; every reference inside the delivered unit resolves inside it, so no unresolved reference or citation is produced. Printed slips kept literal: no printed word, symbol, hypothesis or number of the counted statement or of its proof was altered. Layout and class: the article's own class is \texttt{elsarticle} with packages including geometry, amsmath, amssymb, amsfonts, amsthm, mathrsfs, graphicx, booktabs, algorithm, algpseudocode, xcolor, hyperref, url; the wrapper is the lane's pinned \texttt{amsart} editorial wrapper at 10pt on A4, which restates the theorem-like environments the retained text needs and adds the article's own macro definitions that the retained text needs (\texttt{\textbackslash R}, \texttt{\textbackslash diag}, \texttt{\textbackslash tr}), with extra wrapper packages (bm, bbm, dsfont, xspace, cancel, mathtools). The delivered numbering is the unit's own. Assets: no retained span contains a figure environment, a graphics-inclusion command, a PDF-page inclusion, a table or a TikZ picture; no image file is copied into this package, the package declares no asset dependency and no image file is redistributed with it. Licence: version arXiv:2603.01975v1 of this article is distributed under the Creative Commons Attribution 4.0 International licence, as recorded in the arXiv licence metadata for this exact version and shown on the abstract page https://arxiv.org/abs/2603.01975v1. A full rights notice is delivered at licenses/arxiv-2603.01975-CC-BY-4.0.txt. Characters: every retained excerpt is pure ASCII, so the delivered engine, XeLaTeX with its default Latin Modern text font, reports no missing character.
\begin{Definition}[Class-normalized profiles, amplitude matrix, and class-normalized density operator]
\label{def:class_normalized_profiles}
Let $F=(\mathbf f_1\,|\,\mathbf f_2\,|\,\cdots\,|\,\mathbf f_k)\in\R^{d\times k}$ be the class-conditional count matrix, where
$\mathbf f_y\in\R_{\ge 0}^d$ collects the attribute counts for class $y$.
Assume $\mathbf 1^\top \mathbf f_y>0$ for all $1\le y\le k$ and define the class-normalized profiles
\[
\widehat p_y \in \R^d,\qquad
\widehat p_y(i):=\frac{F_{i,y}}{\sum_{j=1}^d F_{j,y}}
\quad (1\le i\le d,\ 1\le y\le k).
\]
Define the corresponding \emph{amplitude profiles} $\widehat\psi_y:=\sqrt{\widehat p_y}\in\R^d$ (entrywise square root) and the
amplitude matrix
\[
\widehat\Psi := [\widehat\psi_1\,\cdots\,\widehat\psi_k]\in\R^{d\times k}.
\]
Finally, define the associated \emph{class-normalized density operator}
\[
\rho_{\mathrm{CN}}
:= \frac{\widehat\Psi\widehat\Psi^\top}{\tr(\widehat\Psi\widehat\Psi^\top)}
\in\R^{d\times d}.
\]
Equivalently, since $\|\widehat\psi_y\|_2=1$ for each $y$, one has
\[
\tr(\widehat\Psi\widehat\Psi^\top)=k,
\qquad\text{hence}\qquad
\rho_{\mathrm{CN}}=\frac{1}{k}\widehat\Psi\widehat\Psi^\top.
\]
In later probabilistic statements (population/empirical perturbation analysis), we reserve the notation
$\widehat\rho_\star$ for the empirical estimator of the population class-normalized operator $\rho_\star$.
\end{Definition}

\begin{Notation}[Class masses and normalized profiles]
\label{not:class_masses_profiles}
Let $F\in\R^{d\times k}$ be the class-conditional count matrix with columns $\mathbf f_y$.
Define the class mass (total count) of class $y$ by
\[
s_y := \sum_{i=1}^d F_{i,y} = \mathbf 1^\top \mathbf f_y,
\qquad
S := \sum_{y=1}^k s_y.
\]
Assume $s_y>0$ for all $y$.
Define the class-normalized categorical profiles
\[
p_y := \frac{\mathbf f_y}{s_y}\in\Delta^{d-1},
\qquad
\psi_y := \sqrt{p_y}\in\R^d,
\qquad
\Psi := [\psi_1\,\cdots\,\psi_k]\in\R^{d\times k},
\]
where the square root is taken entrywise.
\end{Notation}

\begin{Proposition}[Exact decomposition and imbalance weights]
\label{prop:exact_relation_rho_hat_rho}
With the notation of Definition~\ref{def:class_normalized_profiles} and Notation~\ref{not:class_masses_profiles}, the two operators admit the explicit forms
\[
\rho_{\mathcal D} \;=\; \frac{1}{S}\,\Psi\,\diag(s_1,\dots,s_k)\,\Psi^\top
\;=\;
\sum_{y=1}^k w_y\,\psi_y\psi_y^\top,
\qquad
w_y := \frac{s_y}{S},
\]
and
\[
\rho_{\mathrm{CN}} \;=\; \frac{1}{k}\,\Psi\Psi^\top
\;=\;
\frac{1}{k}\sum_{y=1}^k \psi_y\psi_y^\top.
\]
In particular, $\rho_{\mathcal D}$ is a convex combination of rank-one matrices $\psi_y\psi_y^\top$ weighted by the
empirical class masses $w_y$, whereas $\rho_{\mathrm{CN}}$ assigns uniform weights $1/k$.
\end{Proposition}

\begin{proof}
Since $\mathbf f_y = s_y p_y$, we have entrywise
\[
\sqrt{\mathbf f_y} = \sqrt{s_y}\,\sqrt{p_y} = \sqrt{s_y}\,\psi_y.
\]
Hence
\[
X = [\sqrt{\mathbf f_1}\,\cdots\,\sqrt{\mathbf f_k}]
   = [\sqrt{s_1}\psi_1\,\cdots\,\sqrt{s_k}\psi_k]
   = \Psi\,\diag(\sqrt{s_1},\dots,\sqrt{s_k}).
\]
Therefore
\[
XX^\top = \Psi\,\diag(s_1,\dots,s_k)\,\Psi^\top.
\]
Moreover, $\|\psi_y\|_2^2=\sum_{i=1}^d p_y(i)=1$, so
\[
\tr(XX^\top)=\sum_{y=1}^k s_y\,\|\psi_y\|_2^2=\sum_{y=1}^k s_y = S.
\]
This yields $\rho_{\mathcal D}=\frac{1}{S}\Psi\diag(s_1,\dots,s_k)\Psi^\top=\sum_y (s_y/S)\psi_y\psi_y^\top$.

For the class-normalized operator, note that $\tr(\Psi\Psi^\top)=\sum_{y=1}^k\|\psi_y\|_2^2=k$, hence
\[
\rho_{\mathrm{CN}} = \frac{1}{k}\Psi\Psi^\top = \frac{1}{k}\sum_y \psi_y\psi_y^\top.
\]
\end{proof}

\end{document}
